\runningheader{Computer Science Principles}{AP Practice Exam (Fundamentals)}{Page \thepage\ of \numpages}

\begingradingrange{cspfundamentals}

\question[5]
Convert \(100001_{2}\) to a decimal number.
\begin{choices}
  \choice 29
  \CorrectChoice 33
  \choice 31
  \choice 63
\end{choices}
\begin{solution}
  \(100001_{2} = 32 + 1 = 33\).
\end{solution}

\question[5]
Convert \(111111_{2}\) to a decimal number.
\begin{choices}
  \choice 29
  \choice 33
  \choice 31
  \CorrectChoice 63
\end{choices}
\begin{solution}
  \(111111_{2} = 32 + 16 + 8 + 4 + 2 + 1 = 63\).
\end{solution}

\question[5]
Convert \(2D_{16}\) to a decimal number.
\begin{choices}
  \choice 29
  \choice 33
  \choice 31
  \CorrectChoice 45
\end{choices}
\begin{solution}
  \(2D_{16} = 2 \times 16 + 13 = 45\).
\end{solution}

\question[5]
Convert \(1F_{16}\) to a decimal number.
\begin{choices}
  \choice 29
  \choice 33
  \CorrectChoice 31
  \choice 45
\end{choices}
\begin{solution}
  \(1F_{16} = 1 \times 16 + 15 = 31\).
\end{solution}

\newpage 

\question[5]
Convert \(1101_{2}\) to a decimal number.
\begin{choices}
  \choice 11
  \choice 1
  \CorrectChoice 13
  \choice 1101
\end{choices}
\begin{solution}
  \(1101_{2} = 8 + 4 + 1 = 13\).
\end{solution}

\question[5]
Convert \(48_{16}\) to a decimal number.
\begin{choices}
  \choice 44
  \choice 54
  \CorrectChoice 72
  \choice 89
\end{choices}
\begin{solution}
  \(48_{16} = 4 \times 16 + 8 = 72\).
\end{solution}

\question[5]
Convert \(8_{16}\) to a decimal number.
\begin{choices}
  \choice 44
  \choice 54
  \choice 72
  \CorrectChoice 8
\end{choices}
\begin{solution}
  \(8_{16} = 8\).
\end{solution}

\question[5]
Convert \(AD_{16}\) to a decimal number.
\begin{choices}
  \choice 44
  \choice 54
  \choice 87
  \CorrectChoice 173
\end{choices}
\begin{solution}
  \(AD_{16} = 10 \times 16 + 13 = 173\).
\end{solution}

\newpage

\question[5]
Convert \(AA_{16}\) to a decimal number.
\begin{choices}
  \CorrectChoice 170
  \choice 180
  \choice 190
  \choice 200
\end{choices}
\begin{solution}
  \(AA_{16} = 10 \times 16 + 10 = 170\).
\end{solution}

\question[5]
Order the numbers from lowest to highest.
\begin{enumerate}
  \renewcommand{\labelenumi}{\Roman{enumi}.}
  \item \(12_{16}\)
  \item \(12_{10}\)
  \item \(12_{8}\)
\end{enumerate}
\begin{choices}
  \choice I, II, III
  \choice III, II, I
  \CorrectChoice II, III, I
  \choice I, III, II
\end{choices}
\begin{solution}
  Values are \(18\), \(12\), and \(10\), so II \(<\) III \(<\) I.
\end{solution}

\question[5]
Order the numbers from lowest to highest.
\begin{enumerate}
  \renewcommand{\labelenumi}{\Roman{enumi}.}
  \item \(A_{16}\)
  \item \(11_{10}\)
  \item \(1111_{2}\)
\end{enumerate}
\begin{choices}
  \choice I, II, III
  \choice III, II, I
  \CorrectChoice I, II, III
  \choice I, III, II
\end{choices}
\begin{solution}
  Values are \(10\), \(11\), and \(15\), so I \(<\) II \(<\) III.
\end{solution}

\newpage 

\question[5]
Order the numbers from highest to lowest.
\begin{enumerate}
  \renewcommand{\labelenumi}{\Roman{enumi}.}
  \item \(AA_{16}\)
  \item \(99_{10}\)
  \item \(1111_{2}\)
\end{enumerate}
\begin{choices}
  \choice I, II, III
  \choice III, II, I
  \choice II, III, I
  \CorrectChoice I, II, III
\end{choices}
\begin{solution}
  Values are \(170\), \(99\), and \(15\), so I \(>\) II \(>\) III.
\end{solution}

\question[5]
Order the numbers from highest to lowest.
\begin{enumerate}
  \renewcommand{\labelenumi}{\Roman{enumi}.}
  \item \(AA_{16}\)
  \item \(99_{10}\)
  \item \(77_{8}\)
\end{enumerate}
\begin{choices}
  \choice I, II, III
  \choice III, II, I
  \CorrectChoice I, II, III
  \choice I, III, II
\end{choices}
\begin{solution}
  Values are \(170\), \(99\), and \(63\), so I \(>\) II \(>\) III.
\end{solution}

\question[5]
Convert \(48_{10}\) to a binary number.
\begin{choices}
  \choice \(110001_{2}\)
  \CorrectChoice \(110000_{2}\)
  \choice \(110111_{2}\)
  \choice \(111111_{2}\)
\end{choices}
\begin{solution}
  \(48 = 32 + 16\), so the binary form is \(110000_{2}\).
\end{solution}

\question[5]
Convert \(63_{10}\) to a binary number.
\begin{choices}
  \choice \(110001_{2}\)
  \choice \(110000_{2}\)
  \choice \(110111_{2}\)
  \CorrectChoice \(111111_{2}\)
\end{choices}
\begin{solution}
  \(63 = 32 + 16 + 8 + 4 + 2 + 1\).
\end{solution}

\newpage 

\question[5]
Convert \(17_{8}\) to a binary number.
\begin{choices}
  \choice \(10001_{2}\)
  \CorrectChoice \(1111_{2}\)
  \choice \(10111_{2}\)
  \choice \(11111_{2}\)
\end{choices}
\begin{solution}
  \(17_{8} = 1 \times 8 + 7 = 15\), and \(15 = 1111_{2}\).
\end{solution}


\question[5]
Convert \(AC_{16}\) to a binary number.
\begin{choices}
  \CorrectChoice \(10101100_{2}\)
  \choice \(11001010_{2}\)
  \choice \(11110000_{2}\)
  \choice \(11111111_{2}\)
\end{choices}
\begin{solution}
  \texttt{A} maps to \texttt{1010} and \texttt{C} maps to \texttt{1100}.
\end{solution}

\question[5]
Convert \(2A_{16}\) to a binary number.
\begin{choices}
  \CorrectChoice \(00101010_{2}\)
  \choice \(11001010_{2}\)
  \choice \(11110000_{2}\)
  \choice \(11111111_{2}\)
\end{choices}
\begin{solution}
  \texttt{2} maps to \texttt{0010} and \texttt{A} maps to \texttt{1010}.
\end{solution}

\question[5]
Convert \(16_{16}\) to a binary number.
\begin{choices}
  \choice \(10101_{2}\)
  \CorrectChoice \(10110_{2}\)
  \choice \(11110_{2}\)
  \choice \(11111_{2}\)
\end{choices}
\begin{solution}
  \(16_{16} = 22\), and \(22 = 10110_{2}\).
\end{solution}

\newpage 

\question[5]
Using only \texttt{RANDOM(a, b)} and \texttt{DISPLAY(expression)}, which value is \textbf{not} possible from
\begin{lstlisting}[style=pseudo]
DISPLAY(RANDOM(2, 5))
\end{lstlisting}
\begin{choices}
  \choice 2
  \choice 4
  \choice 5
  \CorrectChoice 6
\end{choices}
\begin{solution}
  \texttt{RANDOM(2, 5)} can only return 2, 3, 4, or 5.
\end{solution}

\question[5]
Using only \texttt{RANDOM(a, b)} and \texttt{DISPLAY(expression)}, which value is \textbf{not} possible from
\begin{lstlisting}[style=pseudo]
DISPLAY(RANDOM(1, 4) + RANDOM(2, 5))
\end{lstlisting}
\begin{choices}
  \choice 6
  \choice 9
  \CorrectChoice 1
  \choice 5
\end{choices}
\begin{solution}
  Minimum possible sum is \(1 + 2 = 3\), so 1 cannot occur.
\end{solution}

\question[5]
Which value will result in overflow in a 3-bit unsigned system?
\begin{choices}
  \choice 3
  \choice 4
  \choice 6
  \CorrectChoice 8
\end{choices}
\begin{solution}
  A 3-bit unsigned range is 0 to 7.
\end{solution}

\question[5]
Which value will result in overflow in a 4-bit unsigned system?
\begin{choices}
  \choice 6
  \choice 9
  \choice 15
  \CorrectChoice 16
\end{choices}
\begin{solution}
  A 4-bit unsigned range is 0 to 15.
\end{solution}

\newpage 

\question[5]
Which is the largest number that will \textbf{not} overflow in a 5-bit unsigned system?
\begin{choices}
  \choice 30
  \choice 34
  \CorrectChoice 31
  \choice 33
\end{choices}
\begin{solution}
  A 5-bit unsigned range is 0 to 31.
\end{solution}


\question[5]
Which expression could result in a round-off error?
\begin{choices}
  \CorrectChoice \(1/3\)
  \choice \(2 \times 6\)
  \choice \(6 - 2\)
  \choice \(3 + 5\)
\end{choices}
\begin{solution}
  \(1/3\) is repeating in binary and may be approximated.
\end{solution}

\question[5]
In an 8-bit system, which hexadecimal number is not representable?
\begin{choices}
  \choice \(14_{16}\)
  \choice \(34_{16}\)
  \choice \(84_{16}\)
  \CorrectChoice \(100_{16}\)
\end{choices}
\begin{solution}
  8 bits can represent \(00_{16}\) to \(FF_{16}\); \(100_{16}\) needs more than 8 bits.
\end{solution}

\question[5]
Why do many systems use fixed-size integers (for example, 4-byte integers)? \textbf{Select two answers.}
\begin{checkboxes}
  \choice Programmers need numbers to be as small as possible because that limits rollover errors.
  \CorrectChoice Most computed numbers are within this range, so 4 bytes is a reasonable size for many purposes.
  \CorrectChoice Making integer size effectively unbounded would require impractical amounts of storage and overhead.
  \choice No programmer ever needs a number outside this range.
\end{checkboxes}
\begin{solution}
  The practical trade-off is memory efficiency and typical use-case coverage.
\end{solution}

\newpage 

\question[5]
A programmer needs to model many properties with the same behavior but different names, areas, and images. What is a benefit of using an abstraction with these as parameters?
\begin{choices}
  \CorrectChoice Parameters let each instance store its own values without duplicating code.
  \choice The abstraction accounts for every possible behavior difference automatically.
  \choice The abstraction primarily secures code from theft.
  \choice The abstraction guarantees each object's files can be edited.
\end{choices}
\begin{solution}
  Parameterization supports reuse and customization per instance.
\end{solution}

\question[5]
Why is debugging low-level machine code generally difficult?
\begin{choices}
  \choice A single bit position can always be identified instantly by inspection.
  \CorrectChoice It is far easier to reason about and debug behavior in a higher-level language.
  \choice Every machine-code bug can be fixed by flipping one specific bit.
  \choice Machine code cannot be tested.
\end{choices}
\begin{solution}
  Higher-level abstractions are easier for humans to understand and verify.
\end{solution}

\question[5]
Order these languages from low-level to high-level.
\begin{enumerate}
  \renewcommand{\labelenumi}{\Roman{enumi}.}
  \item COBOL
  \item Python
  \item Machine code
\end{enumerate}
\begin{choices}
  \choice I, II, III
  \choice II, III, I
  \CorrectChoice III, I, II
  \choice All are at the same abstraction level
\end{choices}
\begin{solution}
  Machine code is lowest-level; Python is highest-level among these.
\end{solution}

\question[5]
Order these hardware layers from low-level abstraction to high-level abstraction.
\begin{enumerate}
  \renewcommand{\labelenumi}{\Roman{enumi}.}
  \item Video card
  \item Transistor
  \item Computer
\end{enumerate}
\begin{choices}
  \choice I, II, III
  \choice II, III, I
  \CorrectChoice II, I, III
  \choice Hardware has no abstraction levels
\end{choices}
\begin{solution}
  Transistors compose components like video cards, which compose full systems.
\end{solution}

\newpage 

\question[5]
Why is high-level source code usually easier to read than low-level code?
\begin{choices}
  \choice It is written by smarter programmers.
  \choice It is closer to machine instructions.
  \choice It is intended only for experts.
  \CorrectChoice It uses abstractions and descriptive constructs closer to natural language.
\end{choices}
\begin{solution}
  Abstraction hides low-level implementation details.
\end{solution}


\question[5]
Taken as a whole, is a physical computer low-level or high-level?
\begin{choices}
  \choice Low-level, because it executes low-level software
  \CorrectChoice High-level, because it is composed of many lower-level subsystems
  \choice High-level, because it can run high-level languages
  \choice Neither, because it is not software
\end{choices}
\begin{solution}
  In hardware abstraction terms, a complete computer sits above its components.
\end{solution}

\question[5]
Logic gates physically implement Boolean operations such as AND and OR. Are Boolean functions abstractions?
\begin{choices}
  \CorrectChoice Yes, they abstract gate behavior into symbolic operations.
  \choice Yes, because they represent the physical gate hardware itself.
  \choice No, because processors perform all logic directly without gates.
  \choice No, because bit-level concepts cannot be abstractions.
\end{choices}
\begin{solution}
  Boolean expressions are conceptual models of underlying hardware behavior.
\end{solution}

\question[5]
A theme park simulation estimates wait times for a popular ride. Which patron characteristics are most useful? \textbf{Select two answers.}
\begin{checkboxes}
  \CorrectChoice Ride preference (which rides the patron chooses)
  \CorrectChoice Walking preference (how far the patron will walk)
  \choice Food preference
  \choice Ticket type
\end{checkboxes}
\begin{solution}
  Ride selection and movement patterns directly affect queue dynamics.
\end{solution}

\newpage 

\question[5]
A simulation models fox and rabbit populations. Which aspect is least necessary to include?
\begin{choices}
  \CorrectChoice A detailed grass model for rabbit diet
  \choice Rabbit litter-size limits
  \choice Foxes needing rabbits to survive
  \choice Rabbit lifespan limits
\end{choices}
\begin{solution}
  In a simplified predator-prey model, grass can often be omitted first.
\end{solution}

\question[5]
Why do FAA investigators use simulators instead of real aircraft when replaying crash scenarios?
\begin{choices}
  \choice Simulators are harder to fly, so success proves any pilot can do it.
  \choice Simulators are easier to fly, so failure proves no pilot could succeed.
  \choice Simulators are more expensive to replace than aircraft.
  \CorrectChoice Simulations avoid risking lives if a scenario still results in a crash.
\end{choices}
\begin{solution}
  Simulation supports safer experimentation.
\end{solution}

\question[5]
A scientist simulates ozone-layer changes over time. Which output is most realistic and useful?
\begin{choices}
  \CorrectChoice An approximate time until recovery
  \choice The exact hole size at any instant
  \choice The exact time until recovery
  \choice The exact depth at any instant
\end{choices}
\begin{solution}
  Simulations are typically used for approximate prediction, not exact certainty.
\end{solution}

\question[5]
Why can simulation be useful for studying invasive-species spread?
\begin{choices}
  \choice A single simulated species can perfectly represent all species at once.
  \choice Simulated species can be created quickly.
  \choice Simulation can run faster than real ecological time.
  \CorrectChoice It enables rapid experiments across many environments and timelines.
\end{choices}
\begin{solution}
  Multiple controlled scenarios can be tested quickly and safely.
\end{solution}

\newpage 

\question[5]
Why is a ``run many trials with fixed variables'' feature useful in a simulation? \textbf{Select two answers.}
\begin{checkboxes}
  \choice It guarantees low hardware usage.
  \CorrectChoice It isolates the effect of specific variables.
  \CorrectChoice It quickly generates a large data set.
  \choice It automatically increases model detail.
\end{checkboxes}
\begin{solution}
  Repeated controlled trials improve analysis quality.
\end{solution}

\question[5]
Which language listed has the highest abstraction level?
\begin{choices}
  \choice COBOL
  \choice Machine code
  \CorrectChoice Python/Java
  \choice All of the above
\end{choices}
\begin{solution}
  Modern high-level languages such as Python and Java provide richer abstractions.
\end{solution}

\question[5]
Which has the lowest level of hardware abstraction?
\begin{choices}
  \CorrectChoice Transistors
  \choice Computer chips
  \choice Motherboard
  \choice Computer
\end{choices}
\begin{solution}
  Transistors are more fundamental than chips, boards, and full systems.
\end{solution}

\question[5]
Which number system has the lowest abstraction level in computing?
\begin{choices}
  \CorrectChoice Binary
  \choice Octal
  \choice Decimal
  \choice Hexadecimal
\end{choices}
\begin{solution}
  Binary maps directly to physical on/off states.
\end{solution}

\endgradingrange{cspfundamentals}